%! AP Statistics — Unit 2, Lesson 2.7 — Group Activity — Answer Key
\documentclass[10pt]{article}
\usepackage{apstats-article}
\usepackage{apstats-key}

\begin{document}

\pageheader{Unit 2, Lesson 2.7}{Group Activity --- Answer Key}
\namepartnerperiod

\vspace{0.18in}

% ── SCENARIO ─────────────────────────────────────────────────────────────────
\begin{scenariobox}[Directions]{sky}
\small
Each tier uses a real-world regression context. Compute residuals, construct or interpret a
residual plot, and write a complete AP-style conclusion about model appropriateness. Write
all sentences in context using the variable names.
\end{scenariobox}

\vspace{0.12in}

% ── TIER R ────────────────────────────────────────────────────────────────────
\begin{tcolorbox}[colback=white, colframe=black!40, boxrule=0.8pt, arc=2mm,
    left=4mm, right=4mm, top=3mm, bottom=3mm,
    title={\textbf{Tier R --- Remediate}}, fonttitle=\bfseries, halign title=center]
\small

\textbf{Context: Age and Blood Pressure}

A health clinic collected data on patient age ($x$, years) and systolic blood pressure ($y$,
mmHg). Technology produced the regression equation:
\[
  \widehat{\text{BP}} = 80 + 0.5(\text{age})
\]
The table below shows age, blood pressure, and predicted blood pressure for five patients.

\vspace{8pt}
\renewcommand{\arraystretch}{1.9}
\begin{tabularx}{\linewidth}{@{} >{\centering\arraybackslash}p{2.4cm}
                                    >{\centering\arraybackslash}p{2.4cm}
                                    >{\centering\arraybackslash}p{2.8cm}
                                    X @{}}
\rowcolor{sky}
\textbf{Age ($x$)} & \textbf{BP ($y$)} & \textbf{Predicted ($\hat{y}$)} &
\textbf{Residual ($e = y - \hat{y}$)} \\
\midrule
20 & 90 & 90 & \ans{0} \\
30 & 96 & 95 & \ans{$+1$} \\
40 & 99 & 100 & \ans{$-1$} \\
50 & 106 & 105 & \ans{$+1$} \\
60 & 109 & 110 & \ans{$-1$} \\
\midrule
\multicolumn{3}{r}{\textbf{Sum of residuals}} & \ans{$0$} \\
\end{tabularx}

\vspace{10pt}

\begin{enumerate}[leftmargin=1.6em, itemsep=0.18in]

  \item Compute and record the residual for each patient. Verify that the residuals sum to 0.\\[6pt]
        (Check: sum $=$ \ans{$0$} ✓)

  \item Interpret the residual for the 30-year-old patient in context.\\[8pt]
        \ansline{The residual is $e = 96 - 95 = +1$ mmHg. The actual blood pressure for the
        30-year-old patient is 1 mmHg higher than the value predicted by the regression line.}

  \item A residual plot for these data shows random scatter around $e = 0$ with no obvious
        pattern. Write a complete conclusion about whether a linear model is appropriate for
        predicting blood pressure from age.\\[8pt]
        \ansline{The residual plot shows no apparent pattern (random scatter around zero).
        This suggests that a linear model is an appropriate model for predicting systolic
        blood pressure from patient age.}

\end{enumerate}
\end{tcolorbox}

\vspace{0.12in}

% ── TIER A ────────────────────────────────────────────────────────────────────
\begin{tcolorbox}[colback=white, colframe=black!40, boxrule=0.8pt, arc=2mm,
    left=4mm, right=4mm, top=3mm, bottom=3mm,
    title={\textbf{Tier A --- Approaching Proficiency}}, fonttitle=\bfseries, halign title=center]
\small

\textbf{Context: Hours of TV and GPA}

A researcher recorded daily TV-watching hours ($x$) and GPA ($y$) for five students.
The data ranged from \textbf{1 to 5 hours} of TV per day. Technology output:
\[
  \widehat{\text{GPA}} = 4.05 - 0.35(\text{TV hours})
\]

\vspace{8pt}
\renewcommand{\arraystretch}{1.9}
\begin{tabularx}{\linewidth}{@{} >{\centering\arraybackslash}p{2.6cm}
                                    >{\centering\arraybackslash}p{2.0cm}
                                    >{\centering\arraybackslash}p{2.8cm}
                                    X @{}}
\rowcolor{sky}
\textbf{TV hours ($x$)} & \textbf{GPA ($y$)} & \textbf{Predicted ($\hat{y}$)} &
\textbf{Residual ($e = y - \hat{y}$)} \\
\midrule
1 & 3.90 & \ans{3.70} & \ans{$+0.20$} \\
2 & 3.00 & \ans{3.35} & \ans{$-0.35$} \\
3 & 3.20 & \ans{3.00} & \ans{$+0.20$} \\
4 & 2.50 & \ans{2.65} & \ans{$-0.15$} \\
5 & 2.40 & \ans{2.30} & \ans{$+0.10$} \\
\midrule
\multicolumn{3}{r}{\textbf{Sum}} & \ans{$0$} \\
\end{tabularx}

\vspace{10pt}

\begin{enumerate}[leftmargin=1.6em, itemsep=0.18in]

  \item Compute $\hat{y}$ and $e$ for each student. Verify $\sum e = 0$ (allow rounding to 0.01).

  \item Plot the residuals on the axes below. Draw a dashed reference line at $e = 0$.

        \vspace{0.1in}
        \begin{center}
        \begin{tikzpicture}
          \draw[->] (0,0) -- (7,0) node[right]{\small TV hours ($x$)};
          \draw[->] (0,-1.6) -- (0,1.6) node[above]{\small residual};
          \draw[dashed,gray] (0,0) -- (7,0);
          \foreach \x/\lbl in {1.2/1, 2.4/2, 3.6/3, 4.8/4, 6.0/5}
            \draw (\x,0.06)--(\x,-0.06) node[below]{\tiny\lbl};
          \foreach \y/\lbl in {-1.2/{$-0.4$}, -0.6/{$-0.2$}, 0.6/{$0.2$}, 1.2/{$0.4$}}
            \draw (0.06,\y)--(-0.06,\y) node[left]{\tiny\lbl};
          % key points: x=1→+0.20, x=2→-0.35, x=3→+0.20, x=4→-0.15, x=5→+0.10
          % scale: 3 units per 0.4 GPA → 1 unit = 0.4/3 → 0.20 → 0.6*3/0.4=1.5? Let's use 1 unit = 0.1 GPA (0.3cm/0.1)
          % Actually use 3 per 0.4 → each 0.1 = 0.3 * (3/4) ... let me just use proportional
          % y-axis: 1.2cm = 0.4 GPA → 1 cm = 0.333 GPA → 1 GPA = 3 cm
          % +0.20 → 0.6cm, -0.35 → -1.05cm, +0.20 → 0.6cm, -0.15 → -0.45cm, +0.10 → 0.3cm
          \filldraw[keyred] (1.2, 0.6) circle (2.5pt);
          \filldraw[keyred] (2.4,-1.05) circle (2.5pt);
          \filldraw[keyred] (3.6, 0.6) circle (2.5pt);
          \filldraw[keyred] (4.8,-0.45) circle (2.5pt);
          \filldraw[keyred] (6.0, 0.3) circle (2.5pt);
        \end{tikzpicture}
        \end{center}

  \item Describe the pattern in the residual plot. Write a complete AP-style conclusion about
        whether a linear model is appropriate for predicting GPA from TV-watching hours.\\[8pt]
        \ansline{The residual plot shows no strong apparent pattern (roughly random scatter
        around zero). This suggests that a linear model is an appropriate model for predicting
        GPA from daily TV-watching hours.}

\end{enumerate}
\end{tcolorbox}

\vspace{0.12in}

% ── TIER E ────────────────────────────────────────────────────────────────────
\begin{tcolorbox}[colback=white, colframe=black!40, boxrule=0.8pt, arc=2mm,
    left=4mm, right=4mm, top=3mm, bottom=3mm,
    title={\textbf{Tier E --- Extension}}, fonttitle=\bfseries, halign title=center]
\small

\textbf{Context: Fertilizer and Plant Growth}

A botanist applied different amounts of fertilizer ($x$, grams per week) to seedlings and
measured their height ($y$, cm) after four weeks. Two different regression models were fit.
Their residual plots are described below.

\vspace{8pt}

\textbf{Model 1 --- Linear:} The residual plot shows a clear U-shaped (curved) pattern.
Residuals are large and positive at low and high fertilizer values, and negative in the middle.

\textbf{Model 2 --- Quadratic:} The residual plot shows random scatter around zero with no
obvious pattern.

\vspace{10pt}

\begin{enumerate}[leftmargin=1.6em, itemsep=0.18in]

  \item Write a complete AP-style statement about whether the linear model is an appropriate
        model for predicting plant height from fertilizer amount.\\[8pt]
        \ansline{The residual plot for the linear model shows a clear curved (U-shaped) pattern.
        This suggests that a linear model is NOT an appropriate model for predicting plant
        height from fertilizer amount; a nonlinear (e.g., quadratic) model may provide a better fit.}

  \item Write a complete AP-style statement about the quadratic model.\\[8pt]
        \ansline{The residual plot for the quadratic model shows no apparent pattern (random
        scatter around zero). This suggests that the quadratic model is an appropriate model
        for predicting plant height from fertilizer amount.}

  \item Which model would you recommend and why?\\[8pt]
        \ansline{The quadratic model is recommended. The linear model's residual plot shows
        a clear curved pattern indicating poor fit, while the quadratic model's residual plot
        shows random scatter, indicating the quadratic model captures the data's curvature.}

  \item The largest absolute residual is at $x = 12$: actual = 32 cm, $\hat{y} = 26$ cm.
        Calculate and interpret this residual.\\[8pt]
        \ansline{$e = 32 - 26 = +6$ cm. The actual plant height at 12 grams/week of fertilizer
        is 6 cm greater than the value predicted by the linear regression line.}

\end{enumerate}
\end{tcolorbox}

\begin{teachernote}
Tier A residuals don't sum to exactly 0 with rounded $\hat{y}$ values. Accept answers within
$\pm 0.01$. The intended sum using full precision is 0. Coach students to divide residuals
evenly above and below zero as a qualitative check rather than algebraic precision.
\end{teachernote}

\end{document}
